
@@M0@@ of locally free @@M1@@-modules of ranks @@M2@@) by gluing the free @@M3@@-module @@M4@@ on local trivializing charts by means of the representation @@M5@@. More precisely, the sequence @@M6@@ corresponds to a right @@M7@@-torsor @@M8@@ such that there is a canonical isomorphism
\sgaStableTarget{sga6.e10.eq.6-1-3.occ2}@@M9@@
where @@M10@@ denotes the contracted product, and @@M11@@ acts on @@M12@@ through the @@M13@@-th projection @@M14@@; then every representation
@@M15@@
makes @@M16@@ act on @@M17@@ by @@M18@@-automorphisms, and thus defines a multifunctor @@M19@@ such that
\sgaStableTarget{sga6.e10.eq.6-1-4}@@M20@@
Moreover, for every composite representation @@M21@@, one finds
\sgaStableTarget{sga6.e10.eq.6-1-5}@@M22@@

\sgaStableTarget{sga6.e10.sec.6.2}\subsection*{6.2.}
For every locally free @@M23@@-module @@M24@@ of rank @@M25@@, we shall denote by @@M26@@ the invertible @@M27@@-module @@M28@@. Keeping the notation of 6.1, the functor @@M29@@ is determined, up to canonical isomorphism, by the corresponding representation
\sgaStableTarget{sga6.e10.eq.6-2-1}@@M30@@

Suppose now that @@M31@@ is the ringed topos of sheaves on the ringed site @@M32@@ endowed with the étale topology. Then one can calculate @@M33@@, or equivalently the representation @@M34@@, by means of the following
